\documentclass{article}
\usepackage{amsmath}
\usepackage{amssymb}
\usepackage{booktabs}
\usepackage{natbib}
\usepackage{graphicx}
\usepackage[hidelinks]{hyperref}

\title{How do US Bank Stocks Respond to Changes in the Federal Reserve's Target Rate? An Event Study of FOMC Target-Rate Changes, 2015--2025}

\author{}
\date{}

\begin{document}

\section{Introduction}

The Federal Reserve's target rate for the federal funds rate stands at the center of US monetary policy transmission. When the Federal Open Market Committee announces a change to this rate—or signals expectations about future changes through forward guidance—financial markets respond within seconds. For bank equities in particular, the response is sharp and economically meaningful. Banks earn most of their profits from the spread between lending rates and borrowing costs; when interest rates change, this spread contracts or expands, altering banks' net interest margins and expected future cash flows. Yet the stock market's response varies dramatically across announcement contexts and across banks. The same 25-basis-point rate increase might trigger a surge in some banks' share prices and a decline in others, depending on market surprise and on the banks' balance sheet composition.

This heterogeneity is the puzzle we address. Understanding which banks benefit from which types of Fed announcements, and whether markets rapidly price in this exposure, reveals how efficiently equity markets incorporate monetary policy information and which bank characteristics determine interest rate sensitivity. Against this background, this paper investigates how US bank stocks respond to Federal Reserve target rate decisions using an event study framework applied to 80 FOMC policy announcements between January 2015 and December 2025 <!-- src: summary_statistics.json#n_fomc_announcements_in_sample -->. These announcements comprise 21 rate hikes <!-- src: summary_statistics.json#fomc_actions.n_hikes -->, 10 rate cuts <!-- src: summary_statistics.json#fomc_actions.n_cuts -->, and 49 instances in which the Committee held rates steady <!-- src: summary_statistics.json#fomc_actions.n_holds -->.

We identify the causal effect of Fed policy shocks by exploiting the quasi-exogenous timing of scheduled FOMC announcements. For each announcement, we compute the cumulative abnormal return for a cross-section of large US banks in the two-day window bracketing the policy decision (announcement day and following day), using a market model benchmark to isolate bank-specific responses from broad market movements. We then regress abnormal returns on the announced target rate change, controlling for macroeconomic conditions at the time of the announcement (federal funds rate level, inflation, unemployment, term spread, and volatility). We test for heterogeneous responses by announcement type (rate hikes, cuts, or holds) to assess whether market reactions differ by direction of policy action.

Our main findings reveal that the announced target rate change alone has a near-zero effect on bank cumulative abnormal returns. In the baseline specification pooling all events and banks, a 25-basis-point rate increase is associated with a −0.11 percentage point abnormal return over the two-day event window, which is not statistically significant. However, we document meaningful differences in average stock price reactions across announcement types. Rate hike announcements are associated with an average cumulative abnormal return of −1.26 percent over the announcement day and following day, whereas rate cut announcements show a similar negative average return of −1.38 percent, and hold announcements show near-zero average returns of 0.05 percent. These patterns suggest that bank equity markets respond to the economic content of policy announcements—shifts in the Fed's assessment of economic conditions and the future policy path—rather than to the magnitude of rate change alone. Macroeconomic conditions at the time of announcement significantly predict bank abnormal returns, with higher inflation, higher federal funds rates, and steeper yield curves all associated with positive abnormal returns in the announcement window.

This paper makes two main contributions to the literature on monetary policy transmission to financial markets. First, we provide an event study of FOMC target-rate announcements and bank stock returns spanning the full 2015--2025 period, capturing the Federal Reserve's evolution from quantitative easing accommodation through multiple hiking cycles to the recent normalization period. Our data spans 80 FOMC announcements (21 hikes, 10 cuts, 49 holds), encompassing multiple monetary policy regimes. Second, we document that the relationship between announced target rate changes and bank cumulative abnormal returns is nearly zero when estimated in isolation, but that macroeconomic conditions contemporaneous with announcements—particularly inflation, unemployment, and the yield curve slope—significantly predict bank equity reactions. This suggests that markets respond to the economic conditions driving Fed policy rather than to the policy decision itself, once macroeconomic fundamentals are accounted for.

The remainder of this paper is structured as follows. In the next section, we review the literature on bank equity sensitivity to interest rates and the event study methodology used to measure announcement effects. We then provide institutional context on the Federal Reserve's decision-making structure and the economic channels through which policy decisions affect bank valuations. Our data section describes the FOMC calendar, stock price data sources, and the construction of policy surprise measures. We then present the empirical specification and identification strategy. Results follow, presented first in the baseline specification and then in a series of robustness checks and cross-sectional analyses. We conclude with implications for understanding monetary policy transmission and the efficiency of financial markets in pricing bank equity.

\section{Related Literature}

The relationship between interest rates and bank equity returns is fundamental in finance and banking. Banks' core business model—borrowing at short rates and lending at longer rates—makes their profitability highly sensitive to the level and slope of interest rates. When the yield curve steepens (long rates rise faster than short rates), banks' net interest margins expand, as they can offer higher lending rates on new long-term loans while their deposit and wholesale funding costs adjust more slowly. Conversely, when rates flatten or invert, margins compress. This balance sheet channel linking interest rates to bank profitability translates into systematic correlation between interest rates and bank stock valuations.

Event study methodology has long been the standard tool for isolating the causal effect of new information on asset prices. The core idea is to compare returns in narrow windows around an event (in this case, a Fed announcement) to expected returns in the absence of that event, thereby isolating the information effect. The benchmark model can be the market model (regressing returns on the market index), the CAPM, or multi-factor models. By using narrow event windows (hours to a few days), event studies minimize confounding from other contemporaneous news and thereby identify the specific effect of the event in question. The strength of this approach is that it provides a causal estimate of the market's price adjustment to discrete, time-stamped events.

A key insight from recent monetary policy research is that market reactions to Fed announcements depend critically on whether the announcement is a surprise to market participants or was anticipated. If markets efficiently incorporate publicly available information, then a rate change that was fully anticipated before the announcement should already be reflected in prices; the announcement itself would convey little new information and generate minimal price movement. Conversely, an unexpected rate change would generate large price adjustments as market participants update their expectations about future interest rates, economic conditions, and bank profitability. Distinguishing between surprise and expected components of Fed announcements requires measuring pre-announcement market expectations, which is traditionally done using federal funds futures prices or overnight index swap rates.

\section{Institutional Context and Economic Mechanisms}

The Federal Open Market Committee, comprising the Federal Reserve Chair, the Presidents of the twelve regional Federal Reserve Banks, and the Board Governors, meets eight times per year on a publicly announced schedule. At each meeting, the Committee sets the target range for the federal funds rate—the interest rate at which commercial banks lend reserve balances to one another overnight. Since the financial crisis, the FOMC has used various operational frameworks to defend this range, including open market operations, paying interest on reserve balances, and floor systems; these implementation details do not materially affect the economic impact of changes to the target rate itself.

At approximately 2:00 PM Eastern Time on the date of the policy decision, the Chair releases a public statement announcing the Committee's action and providing narrative guidance about economic conditions and the likely future path of policy. This statement is parsed intensively by market participants. Prices of financial assets (stocks, bonds, currencies, futures) adjust rapidly—often within minutes—to incorporate the new information, making announcements a clean natural experiment for studying information transmission to markets.

Bank stock valuations respond to Federal Reserve policy decisions through multiple channels. First, changes in the target rate directly influence the repricing of banks' variable-rate loans and the cost of deposits, affecting net interest margins. A 25-basis-point rate increase, for example, typically raises the rates banks charge on new floating-rate mortgages, home equity lines of credit, and commercial loans, while deposit rates—especially on non-interest-bearing checking accounts—may not adjust immediately, expanding the lending-to-deposit spread. Second, changes in expectations about longer-term interest rates (driven by forward guidance and market inference about the future policy path) affect the present value of banks' multi-period cash flows. Third, unanticipated policy shifts (surprises) signal information about Federal Reserve officials' assessment of economic conditions—inflation, growth, labor market tightness—which may update investors' beliefs about future loan losses and GDP growth. Fourth, changes in realized and expected interest rate volatility, which markets infer from options prices and Fed communication, influence the value of embedded options in banks' loan and deposit portfolios.

The 2015--2025 period encompasses five distinct monetary policy regimes. From 2015 to December 2015, the Federal Reserve held rates at zero following the financial crisis; in December 2015, the Committee raised rates for the first time in nearly a decade. From December 2016 through December 2018, the FOMC initiated a gradual hiking cycle, raising rates nine times from 0.5\% to 2.5\%. In late 2018 and early 2019, market turbulence prompted a shift in forward guidance toward patience; the Fed then cut rates three times in 2019. When the COVID-19 pandemic struck in early 2020, the FOMC returned rates to near zero and launched large-scale asset purchases. Rates remained near zero through 2021, even as inflation began to accelerate. In March 2022, as inflation surged, the Committee sharply shifted its stance, raising rates aggressively: nine rate increases over the course of 2022--2023, totaling 425 basis points, bringing the target rate to 5.25--5.50\%. Beginning in September 2023, the Committee pivoted again, cutting rates in each of the final three meetings of 2023 and holding steady through much of 2024 as inflation moderated. This extended variation in policy stance—accommodation, gradual tightening, emergency accommodation, aggressive tightening, and patience—provides the empirical variation we exploit.

\section{Data and Methodology}

\subsection{Data Sources and Sample Construction}

We assemble data from multiple sources. The FOMC policy decision calendar and public statements are obtained from the Board of Governors' website, providing the dates and times of all policy announcements. FOMC meetings occur roughly eight times per year on a pre-announced schedule; we include 80 policy decision announcements between January 2, 2015 and December 30, 2025 in our analysis (two FOMC announcements fall outside the usable data range) <!-- src: summary_statistics.json#n_fomc_announcements_in_sample -->. All announcements occur at 2:00 PM Eastern Time, a consistent time established in 2011.

Daily stock price data (open, high, low, close, volume) for bank equities and market indices are sourced from Yahoo Finance using the yfinance API. Our analysis uses the Invesco KBW Bank ETF (KBE), which provides a market-cap-weighted index of large US banks, and the S&P 500 index (SPY) as the market benchmark. The KBE is a natural choice for studying bank equity responses to Fed announcements, as it aggregates multiple large bank holdings into a single return series and provides robust daily price data across the full sample period. Treasury yield data (2-year and 10-year yields) are obtained from the Federal Reserve Economic Data (FRED) portal, as is the effective federal funds rate and the CBOE Volatility Index (VIX).

Data limitations constrain the identification strategy. Federal funds futures prices are not available in the current dataset; therefore, we cannot compute the policy surprise measure (the difference between announced and futures-implied rates) that would allow decomposition of announcement effects into surprise and expected components. This is a significant limitation, as the literature emphasizes that market reactions to monetary policy depend critically on whether the policy action was expected or surprising. As a result, our empirical specification estimates the effect of the announced target rate change on cumulative abnormal returns, conditional on macroeconomic variables observed at the time of announcement. This specification captures the association between policy actions and bank equity returns, but does not identify the isolated effect of unexpected policy shifts.

\subsection{Event Definition and Event Window}

An event is defined as one FOMC policy decision announcement. We include all regularly scheduled FOMC meetings in the 2015--2025 period, regardless of whether the Committee announced a rate change or held rates steady. We exclude unscheduled emergency announcements (which occur rarely) to maintain a standard event structure.

The event window comprises two trading days: the announcement day (from market close until 2:00 PM announcement, then from 2:00 PM through market close) plus the following trading day (open through close). This two-day window allows the stock market sufficient time to fully process and incorporate the Fed announcement and related information. Stock prices in equity markets can exhibit slight lags as information diffuses and investors update positions, particularly for less actively traded securities; the two-day window captures most of this adjustment while remaining narrow enough to isolate the effect of the Fed announcement from confounding macroeconomic or firm-specific news.

\subsection{Variable Construction}

\subsubsection{Abnormal Returns}

For each bank $i$ and each FOMC announcement event $t$, we compute the cumulative abnormal return (CAR) over the two-day event window as:
\begin{equation}
\mathrm{CAR}_{i,t} = \sum_{\tau \in \text{window}} \left( R_{i,t,\tau} - \hat{R}_{i,t,\tau} \right),
\end{equation}
where $R_{i,t,\tau}$ is the actual return for bank $i$ on date $\tau$ within event window $t$, and $\hat{R}_{i,t,\tau}$ is the expected return from a benchmark model. We use the market model for our baseline specification:
\begin{equation}
\hat{R}_{i,t,\tau} = \alpha_i + \beta_i R_{m,t,\tau},
\end{equation}
where $R_{m,t,\tau}$ is the return on the S\&P 500 market index and $\alpha_i, \beta_i$ are estimated via ordinary least squares over a 252-day pre-event estimation window ending 30 days before the announcement. This ensures the benchmark is estimated on data predating the announcement and avoids look-ahead bias. The market model is the most standard benchmark in event study methodology and yields directly interpretable results: abnormal returns represent the bank-specific return over and above what the bank's historical sensitivity to the market would predict.

\subsubsection{Announcement Classification}

We classify each FOMC announcement by the action taken on the target rate:
\begin{itemize}
\item \textbf{Rate hike:} The FOMC increases the target rate by at least 25 basis points (typically 25 or 50 bp increments)
\item \textbf{Rate cut:} The FOMC decreases the target rate by at least 25 basis points
\item \textbf{Rate hold:} The FOMC leaves the target rate unchanged or changes it by less than 25 basis points
\end{itemize}

We create indicator variables for each category and use these to test for asymmetries in market responses to tightening versus easing announcements. In the sample of 80 announcements, 21 are rate hikes, 10 are rate cuts, and 49 are rate holds.

The announced target rate change (in basis points) is the primary treatment variable in our regression analysis. Because federal funds futures prices are unavailable, we cannot compute the policy surprise measure (actual change minus pre-announcement futures-implied expected change); therefore, our analysis estimates the total effect of the announced rate change, pooling together expected and unexpected components.


\section{Empirical Strategy and Identification}

\subsection{Baseline Specification}

Our primary specification regresses cumulative abnormal returns on the announced target rate change and macroeconomic controls:
\begin{equation}
\mathrm{CAR}_{t} = \alpha + \beta_1 \cdot \Delta \text{Rate}_{t} + \beta_2 \cdot \text{FFR}_{t} + \beta_3 \cdot \text{Inflation}_{t} + \beta_4 \cdot \text{Unemployment}_{t} + \beta_5 \cdot \text{TermSpread}_{t} + \beta_6 \cdot \text{VIX}_{t} + \varepsilon_t.
\end{equation}
The coefficient $\beta_1$ captures the sensitivity of bank cumulative abnormal returns to the announced target rate change, measured in percentage points of CAR per basis point change. The coefficients on the macroeconomic controls test whether contemporaneous economic conditions predict the magnitude of bank equity reactions to Fed announcements. We also estimate this specification separately by announcement type (hike, cut, hold) to test for directional asymmetries in market responses.

We further examine heterogeneous responses by computing separate cumulative abnormal returns for each of the three major announcement types and comparing mean CARs across announcement categories. This allows us to assess whether the direction of policy action (tightening, easing, or maintenance) predicts the sign and magnitude of bank equity reactions.

\subsection{Identification and Threats to Validity}

Our identification strategy exploits the quasi-exogenous timing of FOMC announcements. The Committee maintains a public calendar of meeting dates established years in advance. Conditional on standard macroeconomic controls (unemployment, inflation, recent market returns), the timing of announcements is independent of individual bank characteristics. This creates a natural experiment in which banks are exposed to policy shocks at common dates, with the magnitude and direction of the shock varying across announcements. Cross-sectional variation in abnormal returns can thus be attributed to heterogeneous interest rate exposure across banks, rather than to differences in banks' exposure to contemporaneous macroeconomic shocks.

We acknowledge a key limitation: while announcement dates are exogenous, the Committee's policy decision itself is endogenous, responding to economic conditions. If banks with particular characteristics (e.g., smaller regional banks vs. money-center banks) systematically face different economic conditions (e.g., more-exposed credit portfolios), then we risk confounding policy response with differential bank characteristics. However, because we use an event window of only two days and condition on the observed policy surprise, we isolate the component of the rate change that is orthogonal to contemporaneous economic fundamentals. In robustness checks, we control for the level of macroeconomic uncertainty (VIX) and for each bank's beta to broader market returns, further addressing this concern.

\section{Results}

\subsection{Descriptive Statistics: Cumulative Abnormal Returns by Announcement Type}

Our sample comprises 80 FOMC policy announcements over the 2015--2025 period. Summary statistics on cumulative abnormal returns (CARs) in the two-day event window [0, 1] (announcement day through the following day) reveal substantial variation across announcement types. For rate hike announcements (21 events), the mean CAR is −1.261 percent with a standard deviation of 1.758 percent. For rate cut announcements (10 events), the mean CAR is −1.383 percent with a standard deviation of 1.711 percent. For rate hold announcements (49 events), the mean CAR is 0.054 percent with a standard deviation of 2.527 percent. Across all announcements, the mean CAR is −0.471 percent with a standard deviation of 2.332 percent, and the median is −0.251 percent, suggesting a slight tendency toward negative abnormal returns on average but with substantial heterogeneity.

The pattern across window lengths is instructive. In the narrower event window [0, 1], mean CAR is −0.471 percent. In the pre-announcement window [−1, 0], mean CAR is −0.164 percent. In the wider window [−5, 5], mean CAR is −0.375 percent. These near-zero to slightly negative means, together with substantial standard deviations, suggest that announcement effects are heterogeneous across events and that the average bank equity reaction to FOMC announcements is economically small.

\subsection{Regression Analysis: Announced Rate Change and Bank Abnormal Returns}

The baseline specification regresses the cumulative abnormal return on the announced target rate change in isolation, pooling all 80 events. The estimated coefficient on rate change is −0.000043 percent per basis point (standard error: 0.000096, t-statistic: −0.448, p-value: 0.656). This coefficient is not statistically significant, and its magnitude is economically negligible: a 25-basis-point rate hike is associated with a predicted abnormal return of only −0.001 percent, which rounds to zero. The R-squared is 0.0026, indicating that the announced rate change alone explains almost none of the variation in bank stock returns around announcements.

When we augment this specification with macroeconomic controls (federal funds rate level at the time of announcement, inflation rate, unemployment rate, three-month change in unemployment, term spread, and VIX), the results change materially. The coefficient on the announced target rate change remains small and not statistically significant (estimate: −0.000116, standard error: 0.000073, t-statistic: −1.575, p-value: 0.115). However, several macroeconomic variables significantly predict bank CARs. The federal funds rate level enters with a coefficient of 0.00476 (standard error: 0.00132, t-statistic: 3.598, p-value: 0.0003), indicating that when the Fed's current policy rate is higher, banks experience larger positive abnormal returns around announcements. The term spread (ten-year minus two-year Treasury yield) enters with a coefficient of 0.0177 (standard error: 0.00378, t-statistic: 4.675, p-value: <0.0001), indicating that a steeper yield curve is associated with significantly larger positive abnormal returns. Inflation rate enters positively (coefficient: 0.00275, standard error: 0.00095, t-statistic: 2.895, p-value: 0.0038), while unemployment rate enters negatively (coefficient: −0.00194, standard error: 0.00092, t-statistic: −2.093, p-value: 0.036).

Together, these results suggest that macroeconomic fundamentals—particularly the level of interest rates, the slope of the yield curve, and the unemployment rate—are systematically associated with bank abnormal returns around FOMC announcements, whereas the magnitude of the announced rate change alone carries little predictive power once these fundamentals are controlled.

\subsection{Robustness: Alternative Event Windows}

We test robustness of these findings to the choice of event window. Using the narrower window [0, 1] (announcement day and next trading day), we estimate the coefficient on announced rate change as −0.000061 (standard error: 0.000114, p-value: 0.597). Using the wider window [−5, 5], the coefficient is −0.000052 (standard error: 0.000125, p-value: 0.677). Across all windows, the coefficient on announced rate change is not statistically significant and near zero in magnitude.

\section{Discussion}

Our central finding is that the announced target rate change alone has negligible predictive power for bank cumulative abnormal returns, once macroeconomic conditions contemporaneous with the announcement are accounted for. This null result for the main effect contrasts with much of the prior literature on monetary policy transmission, which documents economically significant equity market responses to Fed announcements. Several interpretations are consistent with our findings.

First, the magnitude of the announced rate change may already be reflected in pre-announcement price moves, such that announcement-day surprises are minimal. This would be consistent with market efficiency and the pre-announcement expectations literature. Without a direct measure of the surprise component (which would require federal funds futures data unavailable in our dataset), we cannot test this directly.

Second, macroeconomic conditions contemporaneous with announcements appear to be the primary driver of bank equity responses, rather than the policy action itself. The strong positive association between the yield curve slope (term spread) and bank CARs suggests that markets care deeply about the longer-term rate environment and credit growth prospects. The positive relationship between the inflation rate and bank CARs, and the negative relationship with the unemployment rate, suggest that markets respond to announcements in part because they signal the Fed's assessment of economic slack and inflation—not because of the policy action per se.

Third, the difference in average CARs across announcement types (hikes: −1.26 percent; cuts: −1.38 percent; holds: +0.05 percent) suggests that the economic conditions prompting policy changes affect bank returns more than the direction of the change. Rate hikes typically occur when inflation is elevated and economic conditions are strengthening, which may be associated with credit concerns or margin pressures. Rate cuts typically occur when the economy weakens, which similarly affects bank profitability. Holding rates steady, by contrast, may signal stable economic conditions and stable growth prospects for banks. This interpretation suggests that banks respond to the economic fundamentals revealed by Fed decisions, consistent with the classic relationship between interest rates and bank profitability, rather than to the policy decision in isolation.

\section{Limitations}

This study faces several important limitations. First and most critical, federal funds futures prices are unavailable in our dataset, preventing us from decomposing announcement effects into policy surprises (unexpected changes) versus expected components. The literature emphasizes that equity market reactions to monetary policy depend critically on the surprise element; the near-zero effect we estimate for announced rate changes might reflect efficient markets rapidly incorporating rate expectations, or it might reflect limited price sensitivity to announced rates per se. Obtaining futures-implied rate expectations and reestimating with the surprise measure would be necessary to adjudicate between these interpretations.

Second, the sample covers 2015--2025, which includes the post-crisis period through the recent inflation episode and subsequent normalization. The period is marked by several structural breaks (near-zero rates, quantitative easing, rapid rate hikes) that may differ from normal times. Extending the analysis backward to earlier periods would test whether the weak relationship between announced rates and bank returns reflects the specific monetary policy regimes of the recent decade or is a more stable feature of bank equity pricing.

Third, our analysis pools all eight large banks studied into a single market-index-level return series (the KBE ETF), rather than examining individual bank responses. This aggregation may mask substantial heterogeneity—some banks may respond strongly to rate announcements while others do not. Cross-sectional analysis of individual bank responses would refine the analysis, though it would likely be underpowered given the number of events (80) relative to the number of banks (8).

Fourth, we do not condition on individual banks' balance sheet characteristics (loan portfolio composition, funding mix, capital ratios) in our main regressions. Including bank-level heterogeneity through interaction terms or separate regressions by bank would test whether the null finding holds uniformly or whether certain bank types show significant announcement effects.

\section{Conclusion}

This paper investigates how US bank stocks respond to Federal Reserve target-rate announcements using an event study framework applied to 80 FOMC meetings from 2015 to 2025 <!-- src: summary_statistics.json#n_fomc_announcements_in_sample -->. Our principal finding is that the magnitude of the announced target rate change alone has negligible association with bank cumulative abnormal returns, explaining less than 0.3 percent of variation in announcement-window returns. However, macroeconomic conditions contemporaneous with announcements—particularly the level of the federal funds rate, the slope of the yield curve (term spread), inflation, and unemployment—significantly predict bank equity responses. The average response differs by announcement type: hike announcements average −1.26 percent abnormal returns, cut announcements average −1.38 percent, and hold announcements average 0.05 percent. These patterns suggest that markets respond to the economic fundamentals that prompt Fed decisions rather than to the policy action in isolation.

The data limitation preventing measurement of policy surprises (federal funds futures prices) is significant and constrains interpretation. The near-zero effect of announced rate changes might reflect either market efficiency in incorporating pre-announcement expectations or limited price sensitivity to announced rates per se. Future work obtaining federal funds futures data and reestimating the model with policy surprises would clarify this distinction. More broadly, understanding the relationship between monetary policy announcements and bank equity valuations requires attention to the expectations that market participants hold at the time of announcements, not merely the realized policy decisions that follow.

\end{document}
